% =====================================================================
%                              PART II
% =====================================================================
\begin{frame}[standout]
Part II\\[6pt]
{\normalsize Statistics for one hypothesis:\\ privileged-information PMI and prefix compatibility (idea 003)}
\end{frame}

\section{II.0\quad Notation and Hypotheses}

\begin{frame}{Part II notation (1/3): decomposing the signal}
Teacher network \emph{without} $z$, with the same prompt template as the student:\ \ $q^0_t=\pitea(\cdot\mid x,h_t)$.
\begin{eqbox}
\[
d_t=\underbrace{\log q_t(y_t)-\log q^0_t(y_t)}_{\mathrm{PMI}_t}
+\underbrace{\log q^0_t(y_t)-\log p_t(y_t)}_{\Delta^{\mathrm{snap}}_t}
\]
\end{eqbox}
\begin{itemize}
  \setlength\itemsep{3pt}
  \item $\mathrm{PMI}_t$ (named after pointwise mutual information): how much seeing $z$ raises the teacher's log-probability of $y_t$.
  \item $\Delta^{\mathrm{snap}}_t$ (snapshot gap): the gap between teacher and student parameters.
  \item If $\bar\theta$ is set to $\theta$ at rollout time, $\Delta^{\mathrm{snap}}_t=0$ on the first mini-batch of a step, so $d_t=\mathrm{PMI}_t$ there and $\mathrm{PMI}_t$ costs nothing extra. On later mini-batches $\Delta^{\mathrm{snap}}_t=-\log r_t$ ($r_t$ from Section C).
\end{itemize}
\end{frame}

% ---------------------------------------------------------------------
\begin{frame}{Part II notation (2/3): prefix compatibility}
For any policy $\pi$, the probability of a prefix is $\pi(h_t\mid x)=\prod_{u=1}^{t-1}\pi(y_u\mid x,h_u)$, and likewise with $z$ in the context. By the chain rule:
\begin{eqbox}
\[
S_t=\sum_{u=1}^{t-1}\mathrm{PMI}_u=\log\frac{\pitea(h_t\mid x,z)}{\pitea(h_t\mid x)},\qquad S_1=0
\]
\end{eqbox}
\begin{itemize}
  \setlength\itemsep{3pt}
  \item $S_t>0$: the prefix is more likely when $z$ is visible; it follows $z$'s route.
  \item $S_t<0$: the prefix has left that route, e.g.\ a different method or an error.
  \item $S_t\approx0$: no evidence yet, typical early in the response.
\end{itemize}
\end{frame}

% ---------------------------------------------------------------------
\begin{frame}{Part II notation (3/3): further symbols}
\begin{symtable}
$\mathcal H(p)=-\sum_{a}p(a)\log p(a)$ & entropy of a distribution $p$ over $\mathcal V$\\
$\mathrm p_{1},\ \mathrm p_{5},\ \mathrm p_{10},\ \mathrm p_{50},\ \dots$ & percentiles, defined like $\pct$\\
$\operatorname{Var}$, $\operatorname{Cov}$, $\operatorname{corr}$ & variance, covariance, Pearson correlation over tokens\\
$T_m(p)$;\ \ $\operatorname{rank}_p(a)$ & the $m$ most likely tokens; one plus the count of more likely tokens\\
$\hat y_t=\arg\max_a q_t(a)$;\ \ $\mathrm{eos}$ & teacher's top-1 token at $t$; end-of-sequence token\\
$\mathcal W$ & reflection words: wait, but, hmm, let me check, actually, alternatively\\
$z^+$;\ \ $z^-$ & correct reference ($z^+=z$); a wrong one in the same template\\
$\pi(y\mid x)$,\ \ $\pi(y\mid x,z)$ & $\prod_{t=1}^{L}\pi(y_t\mid x,h_t)$: whole-response probability (same with $z$)\\
\end{symtable}
\end{frame}

% ---------------------------------------------------------------------
\begin{frame}{Part II: hypotheses to test (none tested yet)}
\small
\begin{itemize}
  \setlength\itemsep{1pt}
  \item \textbf{H1 (derivability).} OPSD helps when the student could derive $z$ itself (e.g.\ an image crop). On long reasoning it hurts when the student cannot derive $z$ (e.g.\ a reference solution), more so the more information $z$ carries.
  \item \textbf{H2 (compatibility, not position).} Whether the teacher's guidance at $t$ helps depends on $S_t$, not on the relative position $t/L$.
  \item \textbf{H3 (reflection tokens).} Tokens from $\mathcal W$ get negative $\mathrm{PMI}_t$ where $S_t\ge0$, and near-zero or positive $\mathrm{PMI}_t$ where $S_t<0$.
\end{itemize}
\footnotesize
\setlength{\tabcolsep}{4pt}\renewcommand{\arraystretch}{1.1}
\begin{tabularx}{\linewidth}{@{}>{\bfseries}l Y@{}}
\toprule
Hypothesis & Key statistics\\
\midrule
H1 & $\mathrm{PMI}_t$ by type of $z$ (II.4); reflection-word rate (II.6); accuracy vs.\ generation budget (II.1)\\
H2 & reliability AUROC, incremental value, position-stratified AUROC, help vs.\ compatibility (II.5)\\
H3 & $\mathrm{PMI}_t$ of reflection words, split by the sign of $S_t$ and by verdict (II.4)\\
\bottomrule
\end{tabularx}
\end{frame}

% =====================================================================
\section{II.1\quad Training Health and Task Metrics}

\begin{frame}{II.1 Training health (1/2): is the run stable?}
\begin{stattable}
Loss and gradient norm & Loss; gradient norm before clipping; fraction of steps skipped because of non-finite gradients. & \cFree & Is training stable?\\
Clip fraction, approximate KL & $\epsilon_{\mathrm{clip}}$: PPO clip range. Fraction of tokens with $r_t\notin[1-\epsilon_{\mathrm{clip}},1+\epsilon_{\mathrm{clip}}]$; mean of $\log\piold(y_t\mid x,h_t)-\log\pit(y_t\mid x,h_t)$ (verl: \texttt{pg\_clipfrac}, \texttt{ppo\_kl}). & \cFree & How far one step moves the policy.\\
Student entropy & $\mathcal H(p_t)$: mean, $\mathrm p_{10}$, $\mathrm p_{50}$, $\mathrm p_{90}$, and the fraction of tokens with $\mathcal H(p_t)<0.01$. & \cFree & Entropy collapse?\\
\end{stattable}
\end{frame}

% ---------------------------------------------------------------------
\begin{frame}{II.1 Training health (2/2): outputs and engines}
\begin{stattable}
Length and truncation & $L$: mean, $\mathrm p_{50}$, $\mathrm p_{90}$, max; fraction of responses cut at the generation limit before $\mathrm{eos}$. & \cFree & Does OPSD lengthen or shorten answers? Is the budget enough?\\
Degenerate output & Fraction of responses with a repeating loop, a format violation (e.g.\ no \texttt{\textbackslash boxed\{\}}), or mixed languages. & \cFree & Is the output degenerating?\\
Engine mismatch & $\pi_{\mathrm{inf}}$: probabilities reported by the inference engine while sampling. Max and $\pct$ of $\abs{\log\pi_{\mathrm{inf}}(y_t\mid x,h_t)-\log\piold(y_t\mid x,h_t)}$ (\texttt{ray\_trainer.py:1153}). & \cFree & Hidden off-policyness.\\
\end{stattable}
\end{frame}

% ---------------------------------------------------------------------
\begin{frame}{II.1 Task metrics (1/2): accuracy and exploration}
\begin{stattable}
Online accuracy & Mean of $c_i$ on the training batch, logged even when $c_i$ is not part of the loss. & \cFree & Is the student improving?\\
pass@$K$ curve & $K$: number of attempts. For a problem with $n_{\mathrm{smp}}$ samples, $n_{\mathrm{cor}}$ of them correct: $\text{pass@}K=1-\binom{n_{\mathrm{smp}}-n_{\mathrm{cor}}}{K}\big/\binom{n_{\mathrm{smp}}}{K}$ (Chen et al.\ 2021), averaged over problems, for $K=1,2,4,\dots,64$. & \cEval & Does distillation narrow what the student can find?\\
maj@$K$ & Fraction of problems whose most frequent answer among $K$ samples is correct. & \cEval & Accuracy with voting.\\
Exploration gap & pass@64 minus pass@1. & \cEval & Is diversity being compressed?\\
Difficulty buckets & $\acc_{\theta_0}(x)$: fraction of correct samples for $x$ under $\theta_0$. Buckets $\{0\}$, $(0,\tfrac12]$, $(\tfrac12,1)$, $\{1\}$; accuracy per bucket over training. & \cEval & Which problems gain, which lose?\\
\end{stattable}
\end{frame}

% ---------------------------------------------------------------------
\begin{frame}{II.1 Task metrics (2/2): regressions, ceiling, budget}
\begin{stattable}
Per-problem flips & A problem is solved if at least half of its samples are correct. Between $\theta_0$ and the final checkpoint: counts of unsolved$\to$solved, solved$\to$unsolved, always unsolved, always solved. & \cEval & Are some problems getting worse?\\
Imitation gap & $\acc$ of the teacher answering with $z$ in context (sampling from $\pitea(\cdot\mid x,z)$) minus $\acc$ of the student. & \cEval & How far could the teacher lift the student?\\
Accuracy vs.\ budget & $\acc(\mathcal D_{\mathrm{held}})$ at several generation limits. & \cEval & Does a longer budget hurt? (reported by Rethinking OPSD)\\
Length by verdict & Mean $L_i$ over responses with $c_i=1$ and with $c_i=0$, separately. & \cFree & How length relates to correctness.\\
\end{stattable}
\vspace{2pt}
\gnote{Out-of-domain accuracy is covered by Part I, Section G.}
\end{frame}

% =====================================================================
\section{II.2\quad The Signal $d_t$}

\begin{frame}{II.2 The signal $d_t$ (1/2): distribution}
\begin{stattable}
Shape & Mean $\mu_d$, standard deviation $\sigma_d$, $\mathrm p_{1}$, $\mathrm p_{5}$, $\mathrm p_{50}$, $\mathrm p_{95}$, $\pct$, min, max, kurtosis $\E[(d_t-\mu_d)^4]/\sigma_d^4$. & \cFree & Is the signal heavy-tailed?\\
Sign and tails & $\Pr(d_t>0)$, $\Pr(d_t<-5)$, $\Pr(\abs{d_t}>2)$. & \cFree & Mostly pushing tokens up, or down?\\
Near-zero share & $\Pr(\abs{d_t}<\epsilon_{\mathrm{low}})$, with $\epsilon_{\mathrm{low}}$ from Section B. & \cFree & How many tokens get almost no supervision?\\
Response-level & $\sum_t d_{i,t}$, $\bar d_i$, and the variance of $d_{i,t}$ within a response. & \cFree & How much responses differ.\\
\end{stattable}
\end{frame}

% ---------------------------------------------------------------------
\begin{frame}{II.2 The signal $d_t$ (2/2): where it falls}
\begin{stattable}
Position buckets & Split $t/L$ into 10 equal buckets; mean and variance of $d_t$ per bucket. & \cFree & Position effect; PW-OPSD and LGR reach opposite conclusions.\\
By verdict & Mean $d_{i,t}$ over responses with $c_i=1$ and with $c_i=0$. & \cFree & Does the teacher keep rewriting answers that are already correct?\\
End-of-sequence signal & $\log q_t(\mathrm{eos})-\log p_t(\mathrm{eos})$ at every position. & \cLog & Is a length change caused by the teacher pushing to stop earlier or later?\\
Token categories & Mean $d_t$ for tokens in $\mathcal W$, digits, formula symbols, format tokens, and answer tokens ($t\in\mathcal S_i$). & \cFree & Which kinds of token carry the signal?\\
\end{stattable}
\end{frame}

% =====================================================================
\section{II.3\quad Teacher--Student Distribution Gap}

\begin{frame}{II.3 Distribution gap (1/2): how far apart are $p_t$ and $q_t$?}
\begin{stattable}
Per-token divergences & $\KL(p_t\Vert q_t)$, $\KL(q_t\Vert p_t)$; $\mathrm{JSD}=\frac12\KL(p_t\Vert m_t)+\frac12\KL(q_t\Vert m_t)$ with $m_t=\frac12(p_t+q_t)$; $\mathrm{TV}=\frac12\sum_a\abs{p_t(a)-q_t(a)}$. & \cLog & Size of the gap; how much switching between forward and reverse KL would matter.\\
Top-1 agreement & $\Pr\big[\arg\max_a p_t(a)=\hat y_t\big]$. & \cLog & Do both want the same next token?\\
Teacher's view of the sample & $q_t(y_t)$ and $\operatorname{rank}_{q_t}(y_t)$. & \cFree\newline\cLog & Does the student often write tokens the teacher finds implausible?\\
Top-10 overlap & $\abs{T_{10}(p_t)\cap T_{10}(q_t)}\big/\abs{T_{10}(p_t)\cup T_{10}(q_t)}$. & \cLog & How similar are the candidate sets?\\
\end{stattable}
\end{frame}

% ---------------------------------------------------------------------
\begin{frame}{II.3 Distribution gap (2/2): entropy and approximation error}
\begin{stattable}
Teacher entropy & $\mathcal H(q_t)$, and the 2-D histogram of $(\mathcal H(p_t),\mathcal H(q_t))$. & \cLog & Where the teacher is sure and the student hesitates (SEAD partitions tokens this way).\\
Top-$m$ mass & $\sum_{a\in T_m(q_t)}q_t(a)$ for the $m$ kept by a truncated loss. & \cLog & Mass ignored by a top-$m$ approximation of the KL.\\
Single sample vs.\ exact & $-d_t$ is an unbiased estimate of $\KL(p_t\Vert q_t)$; report $\operatorname{corr}(-d_t,\KL(p_t\Vert q_t))$. & \cLog & Noise of the single-sample signal (vOPD); the gradient-level version is in Section C.\\
\end{stattable}
\end{frame}

% =====================================================================
\section{II.4\quad PMI Decomposition and Prefix Compatibility}

\begin{frame}{II.4 Identity checks: run before anything else}
\small
On about 20 responses, before trusting any statistic in this part:
\begin{enumerate}
  \setlength\itemsep{4pt}
  \item \textbf{Sum identity.} Check
  \[
  \sum_{t=1}^{L}\mathrm{PMI}_t=\log\pitea(y\mid x,z)-\log\pitea(y\mid x)
  \]
  with whole-response log-probabilities computed separately; the error per token should be below $10^{-3}$.
  \item \textbf{Same tokens.} The response token ids are identical in the two prompt templates (with and without $z$).
  \item \textbf{Logged signal.} With $\bar\theta$ synchronized to $\theta$, the $d_t$ logged during training equals $\mathrm{PMI}_t$ recomputed offline.
\end{enumerate}
\gnote{Cost: \cFwd{} once. These complement the general sanity checks of Section A.}
\end{frame}

% ---------------------------------------------------------------------
\begin{frame}{II.4 Per-token features (idea 003, F0--F7)}
\scriptsize
\setlength{\tabcolsep}{4pt}\renewcommand{\arraystretch}{1.12}
\begin{tabularx}{\linewidth}{@{}>{\bfseries}l >{\raggedright\arraybackslash}p{3.3cm} Y >{\centering\arraybackslash}p{1.6cm}@{}}
\toprule
ID & Feature & Definition & Cost\\
\midrule
F0 & relative position (baseline) & $(t-0.5)/L$ & \cFree\\
F1 & teacher entropy (baseline) & $\mathcal H$ of $q_t$ restricted to $T_{16}(q_t)$, renormalized (PW-OPSD) & \cLog\\
F2 & PMI & $\mathrm{PMI}_t$ & \cFree{} / \cFwd\\
F3 & cumulative PMI & $S_t$ & as F2\\
F4 & mean cumulative PMI & $S_t/(t-1)$ for $t\ge2$ & as F2\\
F5 & local privileged divergence & $\KL(q_t\Vert q^0_t)$ on $T_{64}(q_t)\cup T_{64}(q^0_t)$ plus one bucket for the rest & \cLog\\
F6 & candidate PMI & $\log q_t(a)-\log q^0_t(a)$ for $a\in T_3(q_t)$: max and mean & \cLog\\
F7 & contrastive cumulative PMI & $\sum_{u=1}^{t-1}\big[\log\pitea(y_u\mid x,z^+,h_u)-\log\pitea(y_u\mid x,z^-,h_u)\big]$ & \cFwd\\
\bottomrule
\end{tabularx}
\vspace{2pt}
\gnote{F2--F4 are \cFree{} when $\bar\theta=\theta$. F7 separates the content of $z$ from reference style. II.5 tests F0--F7 as reliability predictors.}
\end{frame}

% ---------------------------------------------------------------------
\begin{frame}{II.4 PMI statistics during training (1/2)}
\gnote{Here and on the next slide, \cFree{} applies when $\bar\theta=\theta$, \cFwd{} otherwise.}\par\vspace{1pt}
\begin{stattable}
PMI distribution & Distribution of $\mathrm{PMI}_t$, overall and by position bucket and token category. & \cFree\newline\cFwd & Which tokens does $z$ favor or penalize?\\
Variance decomposition & $\operatorname{Var}(d_t)=\operatorname{Var}(\mathrm{PMI}_t)+\operatorname{Var}(\Delta^{\mathrm{snap}}_t)+2\operatorname{Cov}(\mathrm{PMI}_t,\Delta^{\mathrm{snap}}_t)$. & \cFwd & With a frozen or EMA teacher, which part drives $d_t$?\\
Compatibility profile & Percentiles of $S_t$ against $t$; fraction of tokens with $S_t<0$; distribution of $S_{L+1}$ on correct and on wrong responses. & \cFree\newline\cFwd & Are prefixes following $z$'s route?\\
Compatibility vs.\ signal & $\operatorname{corr}(S_t,\abs{d_t})$. & \cFree\newline\cFwd & Is the signal more extreme on incompatible prefixes?\\
\end{stattable}
\end{frame}

% ---------------------------------------------------------------------
\begin{frame}{II.4 PMI statistics during training (2/2)}
\begin{stattable}
Reflection-word PMI (H3) & Mean $\mathrm{PMI}_t$ where $y_t\in\mathcal W$: correct vs.\ wrong responses; within wrong ones, $S_t\ge0$ vs.\ $S_t<0$. Control: position-matched tokens not in $\mathcal W$. & \cFree\newline\cFwd & Does suppression depend on the prefix state?\\
Contrastive PMI & $\log\pitea(y_t\mid x,z^+,h_t)-\log\pitea(y_t\mid x,z^-,h_t)$. & \cFwd & Content effect without reference style (RLCSD).\\
PMI by type of $z$ & Same statistics with $z$ = final answer only, full solution, or a hint. & \cFwd{} per type & More information, more harm? (H1)\\
Sequence level & With $\bar\theta=\theta$: $\E_{y\sim\pit(\cdot\mid x)}[S_{L+1}]=-\KL\big(\pit(\cdot\mid x)\,\Vert\,\pit(\cdot\mid x,z)\big)$. & \cFree & Is training pulling $\pit(\cdot\mid x)$ toward $\pit(\cdot\mid x,z)$?\\
\end{stattable}
\end{frame}

% =====================================================================
\section{II.5\quad Prefix Reliability: Branch Continuations}

\begin{frame}{II.5 Branch labels: when does following the teacher help?}
\small
\textbf{Spine}: a fixed student response (correct: $c=1$; wrong: $c=0$). $h\oplus a$: prefix $h$ followed by token $a$. $\hat V(h)$: fraction of $n_{\mathrm{roll}}$ student continuations from $h$, without $z$, that are correct.
\begin{columns}[T]
\column{0.49\textwidth}
\begin{block}{Correct spines (PW-OPSD protocol)}
\footnotesize
Up to 5 positions per problem with the highest F1; force each $a\in T_3(q_t)$; $n_{\mathrm{roll}}=6$.\\
\textbf{diversity}: at least two candidates reach $\hat V(h_t\oplus a)\ge0.75$.\\
\textbf{real-uncertain}: all candidates, and their mean, below $0.40$.\\
Otherwise \textbf{gray}, excluded from AUROC.
\end{block}
\column{0.47\textwidth}
\begin{block}{Wrong spines (help label)}
\footnotesize
5 positions per spine, spread evenly over $t/L$; $n_{\mathrm{roll}}=8$:
\[
\Delta_t=\hat V(h_t\oplus\hat y_t)-\hat V(h_t\oplus y_t)
\]
$\Delta_t>0$: the teacher's top-1 token helps a student without $z$. If $\hat y_t=y_t$, set $\Delta_t=0$ and report how often this happens.
\end{block}
\end{columns}
\vspace{2pt}
\gnote{Cost \cBr: problems $\times5\times3\times6$ continuations for correct spines, problems $\times5\times2\times8$ for wrong spines.}
\end{frame}

% ---------------------------------------------------------------------
\begin{frame}{II.5 Statistics on the branch labels}
\begin{stattable}
Reliability AUROC & AUROC of each of F0--F7 for real-uncertain vs.\ diversity, after subtracting each problem's mean score; problem-level bootstrap 95\% CI (2000 resamples). & \cBr & Can compatibility beat position?\\
Incremental value & 5-fold logistic regression, folds split by problem: held-out AUROC of F0 plus one feature minus that of F0 alone, with bootstrap CI. & \cOff & Information beyond position.\\
Position-stratified AUROC & AUROC of F3 and F5 within each third of $t/L$. & \cOff & Separates compatibility from its correlation with position.\\
Help vs.\ compatibility & $\rho_{\mathrm S}(\Delta_t,S_t)$; partial correlation given $t/L$ (of residuals after regressing both on $t/L$); mean $\Delta_t$ for $S_t\ge0$ vs.\ $S_t<0$. & \cBr & Does leaving $z$'s route distort guidance (LGR)?\\
Rescue and endorsement & Rescue: mean $\hat V(h_t\oplus\hat y_t)$ where $\hat V(h_t)\le0.1$. Error step: $\hat V(h_{t+1})\le\hat V(h_t)-0.3$; fraction of error steps with $\mathrm{PMI}_t>0$. & \cBr & Can the teacher rescue a failing student? Does it endorse mistakes?\\
\end{stattable}
\end{frame}

% ---------------------------------------------------------------------
\begin{frame}{II.5 Pre-registered criteria (from idea 003)}
\footnotesize
\setlength{\tabcolsep}{4pt}\renewcommand{\arraystretch}{1.12}
\begin{tabularx}{\linewidth}{@{}Y Y@{}}
\toprule
\textbf{Result} & \textbf{Conclusion}\\
\midrule
AUROC of F3 or F5 $\ge$ that of F0, and the CI of the increment excludes 0 & Supports H2: compatibility adds information beyond position\\
Increment CI includes 0, but F3/F5 AUROC $>0.6$ within position thirds & Compatibility largely overlaps with position; H2 holds only partly\\
F3/F5 AUROC $\approx0.5$ within position thirds & H2 fails; the position effect has another cause\\
Partial correlation of $\Delta_t$ and $S_t$ given $t/L$ clearly positive & Supports ``leaving $z$'s route distorts guidance''\\
$\Delta_t$ unrelated to $S_t$ & The LGR decay is not explained by compatibility\\
Reflection-word PMI clearly higher on $S_t<0$ segments & Supports H3: masking by prefix state is worth trying\\
F7 beats F3 & Style is a source of noise; the contrastive baseline is needed\\
\bottomrule
\end{tabularx}
\gnote{Keep positive and negative results. One model on one set of problems is not a general conclusion.}
\end{frame}

% =====================================================================
\section{II.6\quad Reasoning Behavior and Diversity}

\begin{frame}{II.6 Reasoning behavior and diversity (1/2)}
\begin{stattable}
Reflection-word rate & Occurrences of words from $\mathcal W$ per 1000 tokens. & \cFree & Fork suppression: does the student stop checking itself?\\
Self-correction & Fraction of responses whose stated answer changes mid-response; of those, the fraction changed to a correct answer and to a wrong one. & \cFree & Is error correction being suppressed?\\
Same-problem diversity & Among the $G$ responses: number of distinct final answers; self-BLEU (Zhu et al.\ 2018); mean pairwise embedding distance; mean position of the first token where two responses differ. & \cFree / \cOff & Are responses collapsing to one form?\\
\end{stattable}
\end{frame}

% ---------------------------------------------------------------------
\begin{frame}{II.6 Reasoning behavior and diversity (2/2)}
\begin{stattable}
Solution routes & Cluster the $G$ responses by method with an LLM judge; number of clusters. & \cEval & Is the reference's route the only one left?\\
Answer leakage & Overlap between $y$ and $z$: shared $n$-grams, ROUGE-L, longest common substring; relative position at which the final answer first appears in $y$. & \cFree & Reasoning or reciting? An answer appearing ever earlier is a warning sign.\\
Style drift & Perplexity under the initial model, $\exp\!\big(-\frac1L\sum_t\log\pi_{\theta_0}(y_t\mid x,h_t)\big)$. & \cFwd & How far has the writing moved from the initial model?\\
\end{stattable}
\vspace{2pt}
\gnote{Here $n$ in ``$n$-gram'' is the span length, e.g.\ 4. Hedging frequency and semantic entropy are in Section E.}
\end{frame}

% =====================================================================
\section{II.7\quad Parameter Dynamics and Cost}

\begin{frame}{II.7 Parameter dynamics}
\begin{stattable}
Drift from the initial model & Mean over probe tokens of $\KL\big(\pit(\cdot\mid x,h_t)\Vert\pi_{\theta_0}(\cdot\mid x,h_t)\big)$. & \cFwd\newline\cLog & How far has the model moved?\\
Teacher drift & $q^{(k)}_t$: teacher distribution with the parameters of step $k$. On $\mathcal P$: $\KL(q^{(k)}_t\Vert q^{(k-N)}_t)$. & \cFwd\newline\cLog & How fast is the target moving when $\bar\theta$ follows $\theta$?\\
Distillation vs.\ RL gradients & $g_{\mathrm D}$, $g_{\mathrm{RL}}$: gradients of the distillation and RL terms of a mixed loss. $\norm{g_{\mathrm D}}/\norm{g_{\mathrm{RL}}}$, $\cos(g_{\mathrm D},g_{\mathrm{RL}})$, fraction of layers with negative cosine. & \cGrad & Are the two objectives fighting? (Dr.OPD)\\
Per-layer update size & $\theta^{(\ell)}$: parameters of layer $\ell$; $\norm{\Delta\theta^{(\ell)}_k}/\norm{\theta^{(\ell)}_k}$. & \cFree & Which layers change?\\
Update structure & For each weight matrix $W$, its change $\Delta W$ over a training window: $\operatorname{erank}(\Delta W)$ (Section D) and the fraction of entries with $\abs{\Delta W}$ below a threshold $\delta_W$. & \cOff & Is the update low-rank or sparse?\\
\end{stattable}
\end{frame}

% ---------------------------------------------------------------------
\begin{frame}{II.7 Cost}
\begin{stattable}
Token counts & Per step: tokens generated, tokens scored by the teacher, extra tokens scored to obtain $q^0_t$. & \cFree & What each piece costs.\\
Time & Share of step time spent on teacher scoring; wall-clock per step; GPU-hours to reach a target accuracy. & \cFree & Where the time goes.\\
Truncated supervision & If the loss keeps a subset of tokens and reweights them by inverse inclusion probability (e.g.\ STOPD): the kept fraction and the largest weight. & \cFree & Is the reweighting stable?\\
\end{stattable}
\end{frame}

% =====================================================================
\section{II.8\quad Slicing, Schedule, Staging}

\begin{frame}{II.8 Slicing and schedule}
\small
\textbf{Slice almost every statistic three ways:} by verdict $c_i$; by initial difficulty bucket (II.1); by position bucket of $t/L$ (II.2). Many effects show only inside a slice, e.g.\ the teacher still restyling responses that are already correct.\\[4pt]
\scriptsize
\setlength{\tabcolsep}{4pt}\renewcommand{\arraystretch}{1.2}
\begin{tabularx}{\linewidth}{@{}>{\bfseries}l >{\raggedright\arraybackslash}p{2.4cm} Y@{}}
\toprule
Tier & When / data & Part II statistics\\
\midrule
1 & every step, full training batch & II.1 training health, online accuracy, length by verdict; II.2; II.4 when $\bar\theta=\theta$: percentiles of $S_t$, fraction with $S_t<0$, reflection-word PMI; II.6 reflection-word rate, self-correction, leakage; II.7 per-layer update size and cost.\\
2 & every $N$ steps, probe set $\mathcal P$ & II.3; II.4 statistics that need $q^0_t$ or $z^-$; II.6 diversity and style drift; II.7 drift and teacher drift.\\
3 & offline & II.1 evaluation metrics; II.4 identity checks; II.5 branch labels and their statistics; II.7 update structure.\\
\bottomrule
\end{tabularx}
\end{frame}

% ---------------------------------------------------------------------
\begin{frame}{II.8 What to add to the dump, and staging}
\begin{columns}[T]
\column{0.47\textwidth}
\begin{block}{Extra fields (on top of Part I's dump)}
\footnotesize
Per token: $\log q^0_t(y_t)$, $\mathcal H(p_t)$, $\mathcal H(q_t)$ as 4-byte floats, plus a 1-byte flag for $y_t\in\mathcal W$: 13 bytes.\\[2pt]
Per response: type of $z$; relative position of the first appearance of the final answer.\\[2pt]
Per branch: position $t$, forced token, $n_{\mathrm{roll}}$, $\hat V$.
\end{block}
\column{0.49\textwidth}
\begin{block}{Staging (idea 003)}
\footnotesize
\textbf{A.} 60 problems: identity checks and correct-spine labels; confirm that the position AUROC is close to PW-OPSD's before comparing new features.\\[2pt]
\textbf{B.} Enough problems for at least 40 real-uncertain positions (keep problems whose initial pass rate is in $[0.2,0.8]$); all branch statistics and H3.\\[2pt]
\textbf{C.} Add the tier-1 compatibility statistics to an ongoing OPSD or STOPD run.
\end{block}
\end{columns}
\end{frame}

